% ──────────────────────────────────────────────────────────────────────
% Comparative Study Research Paper (double-blind submission)
% Springer LNCS proceedings format (llncs.cls, Version 2.25 of 2026/09/03)
%
% Compile with:
%   pdflatex paper
%   bibtex   paper
%   pdflatex paper
%   pdflatex paper
% ──────────────────────────────────────────────────────────────────────
\documentclass[a4paper,orivec]{llncs}
%
\usepackage[T1]{fontenc}
% T1 fonts will be used to generate the final print and online PDFs,
% so please use T1 fonts in your manuscript whenever possible.
%
\usepackage{amsmath}
%
% Springer's official LaTeX template is Computer Modern based, and the author
% instructions ask that the font match the template used (CMR for LaTeX).
% Loading newtx/Times here made the output depend on which optional packages a
% given machine happened to have installed, so the text face is left at the
% class default throughout.
%
\usepackage{graphicx}
\usepackage{array}
\usepackage{url}
\emergencystretch=1.5em       % avoid ugly overfull lines in justified text
% Compact rendering of a 95% confidence interval inside result tables.
\newcommand{\ci}[1]{{\scriptsize\,(#1)}}
% Tighten float and display spacing to keep the paper within the page limit
% without changing the LNCS text block, font sizes, or margins.
\setlength{\textfloatsep}{12pt plus 2pt minus 2pt}
\setlength{\floatsep}{8pt plus 2pt minus 2pt}
\setlength{\intextsep}{10pt plus 2pt minus 2pt}
\setlength{\abovecaptionskip}{6pt}
\setlength{\belowcaptionskip}{0pt}
\setlength{\abovedisplayskip}{7pt plus 2pt minus 3pt}
\setlength{\belowdisplayskip}{7pt plus 2pt minus 3pt}
\setlength{\abovedisplayshortskip}{4pt plus 1pt minus 2pt}
\setlength{\belowdisplayshortskip}{5pt plus 2pt minus 3pt}
% Let floats occupy more of a page and text wrap around them, which removes the
% large vertical gaps that otherwise appear when a table cannot fit in the
% remaining space. Margins, font sizes, and the LNCS text block are unchanged.
\renewcommand{\topfraction}{0.92}
\renewcommand{\bottomfraction}{0.75}
\renewcommand{\textfraction}{0.06}
\renewcommand{\floatpagefraction}{0.80}
\renewcommand{\dbltopfraction}{0.92}
\renewcommand{\dblfloatpagefraction}{0.80}
% Clickable cross-references, citation links, and reference links. hidelinks
% keeps the printed appearance unchanged (no coloured boxes), so the manuscript
% still meets LNCS style rules while the PDF becomes navigable.
\usepackage[hidelinks,breaklinks=true]{hyperref}
\hypersetup{pdftitle={A Comparative Study of Wav2Vec~2.0 and Compact CNN Models
for Cross-Corpus Speech Dysfluency Detection},
pdfsubject={Cross-corpus speech dysfluency detection: Wav2Vec~2.0 compared
with compact spectrogram CNNs},
pdfkeywords={speech dysfluency detection, Wav2Vec~2.0, spectrogram CNN,
cross-corpus evaluation, distribution shift}}
%
% If you use the hyperref package, please uncomment the following two lines
% to display URLs in blue roman font according to Springer's eBook style:
%\usepackage{color}
%\renewcommand\UrlFont{\color{blue}\rmfamily}
%\urlstyle{rm}
%
\begin{document}
%
\title{A Comparative Study of {Wav2Vec}~2.0 and Compact {CNN} Models for
Cross-Corpus Speech Dysfluency Detection}
%
\titlerunning{Wav2Vec~2.0 vs. Compact CNNs}
% If the paper title is too long for the running head, you can set
% an abbreviated paper title here.
%
\author{}
%
\authorrunning{}
% Double-blind submission: author names, affiliations, and contact details
% are deliberately omitted for anonymous review. Restore the real \author{}
% and \institute{} blocks when producing the camera-ready version.
%
\institute{}
%
\maketitle              % typeset the header of the contribution
%
\begin{abstract}
We compare two model families for cross-corpus speech dysfluency detection:
self-supervised Wav2Vec~2.0 and compact mel-spectrogram CNNs. F1 depends on the
decision threshold as well as the model, and thresholds calibrated on one corpus
need not transfer. Seven configurations are compared: five independent binary
Wav2Vec2 classifiers (arm01); two shared-backbone multitask variants frozen for 3
and 20 epochs (arm02, arm03); and four compact mel-spectrogram CNNs with
different sequence aggregators (arm04--arm07). All are trained on 29{,}296
clips from SEP-28K and UCLASS, then evaluated under a
clip-level in-distribution split that shares speakers and a speaker-independent
cross-corpus split on 53 English Project Boli recordings. In-distribution, arm01 attains
the highest macro F1 (0.570, 95\% CI 0.554--0.587, validation-tuned
thresholds), with arm02 close behind (0.556, 0.540--0.572) at
${\approx}4.9\times$ fewer parameters; the intervals overlap. Cross-corpus, the
ordering inverts: an all-positive classifier attains macro F1 $0.550$
(0.490--0.604), which nothing exceeds, because the CNNs reach their F1
by labelling nearly every external clip positive (macro predicted-positive rate
0.879, 95\% CI 0.849--0.909, against a base rate of 0.408) while the Wav2Vec2
configurations collapse toward all-negative (0.079--0.102). Threshold-free
metrics reverse the ranking again: the task-adapted Wav2Vec2 configurations
rank external clips best (macro AUROC 0.620--0.621, AUPRC 0.546--0.557) and
the CNNs lie near chance. The predicted-positive rate shows the CNN operating
points already sat far above the in-distribution base rate, so the CNN advantage
under thresholded F1 is consistent with operating-point drift and does not
imply superior ranking. Thresholded F1 alone therefore cannot support
cross-corpus claims. Each configuration was trained once;
intervals reflect clip-level resampling, and the external analysis is
exploratory.
\end{abstract}
%
\keywords{Speech dysfluency detection \and Wav2Vec 2.0 \and spectrogram CNNs
\and cross-corpus evaluation \and distribution shift \and speaker-independent
evaluation.}
\section{Introduction}
\label{sec:intro}
Stuttering affects approximately 1\% of the global adult population and up to
5\% of children~\cite{yairi2009childhood}. It is characterized by involuntary
repetitions, abnormal prolongations, and pauses (blocks), and it affects
psychosocial well-being. Clinical assessment relies on subjective listening and
waveform inspection, which is time-consuming, subject to inter-rater variability,
and does not scale; this has motivated substantial work on automatic stuttering
identification~\cite{sheikh2022review}. Self-supervised speech representation
learning, and Wav2Vec~2.0 in particular~\cite{baevski2020wav2vec}, has since
produced strong downstream results from raw audio. Applied to dysfluency
detection, however, this work is dominated by a single evaluation style: models
are trained and tested on a common corpus, and performance is summarized with
per-class and macro F1 at a threshold tuned on held-out data from that same
corpus. Such F1 measures two things at once: the quality of the model's ranking of
clips, and the suitability of a threshold calibrated on data resembling the test
data. Only the first should transfer to new speakers and corpora. The two can move
in opposite directions under shift: a threshold that maximizes validation F1
in-distribution may fall in a degenerate region of the ROC curve once the score
distribution changes, leaving F1 dominated by class base rates rather than
discrimination. The available external corpora are small enough that a single
operating point can dominate the reported result.

Cross-corpus and speaker-independent evaluation is not new in principle. Bayerl
et al.~\cite{bayerl2023crosscorpus} studied cross-corpus stuttering detection as
a multi-label problem, and~\cite{bayerl2022partitioning} showed that the choice of
dataset partitioning materially changes reported results. The interaction between
the two metric families under shift is less well characterized: how much of a
cross-corpus F1 difference is attributable to ranking quality, and how much to
operating-point drift. Without that separation, comparisons between representation
families risk being read as statements about generalization when they concern
threshold placement. We address this with a controlled comparison of seven
configurations, Wav2Vec~2.0 models and compact spectrogram CNNs, under matched
protocols, reporting thresholded and threshold-free views side by side. Our
research questions are as follows.
\begin{itemize}
\item \textbf{RQ1}: Under clip-level in-distribution evaluation, how do
Wav2Vec2-based and compact spectrogram CNNs compare in macro and per-class F1?
\item \textbf{RQ2}: On a speaker-independent external corpus, how do their
measured F1 scores and their ordering change?
\item \textbf{RQ3}: How much of the cross-corpus F1 behaviour is explained by
operating-point drift (a mismatch between predicted-positive rate and base rate)
rather than by a change in ranking quality?
\item \textbf{RQ4}: Do threshold-free metrics (AUROC, AUPRC) give a materially
different account than thresholded F1?
\end{itemize}

No configuration generalizes across corpora reliably
(Section~\ref{sec:results}). Our contributions are: (i) a matched comparison of
Wav2Vec~2.0 and compact spectrogram CNNs under in-distribution and
speaker-independent cross-corpus protocols; (ii) an explicit separation of ranking
quality from operating-point placement, using the predicted-positive rate measured
on both splits under identical thresholds, which shows the two metric families
yielding opposite orderings externally; (iii) the observation that no
configuration exceeds an all-positive baseline on external F1, a negative result;
and (iv) a per-class analysis identifying interjection as easiest and block as
hardest.

% ── II. RELATED WORK ─────────────────────────────────────────────────
\section{Related Work}
\label{sec:related}

\subsection{Detection Approaches}
Early systems used handcrafted features (jitter, shimmer, spectral flux, formant
trajectories) with SVMs, HMMs, or GMMs~\cite{sheikh2022review} and generalized
poorly across speakers, accents, and conditions. End-to-end deep learning removed
the manual feature-engineering step: CNNs on mel-spectrograms with recurrent
layers showed early promise for multi-type dysfluency
detection~\cite{kourkounakis2020detecting}, as did RNNs and time-delay networks
for temporal dependencies~\cite{sheikh2021stutternet}. These remain attractive
where labelled data and compute are scarce, and define the compact end of the
capacity range compared against here. Self-supervised pre-training has since
become dominant: Wav2Vec~2.0 contrasts quantized latent speech units against
masked waveform spans and generalizes across languages when pre-trained
multilingually (XLSR-53~\cite{conneau2021xlsr}), HuBERT~\cite{hsu2021hubert}
predicts masked hidden units from offline clustering, XLS-R~\cite{babu2022xlsr}
scales the objective to 128 languages, and WavLM~\cite{chen2022wavlm} adds
denoising objectives. Bayerl et al.~\cite{bayerl2022multi} applied Wav2Vec2 to
multi-task stuttering detection on FluencyBank (macro F1 0.56--0.63 across five
types) and to UCLASS therapy sessions~\cite{bayerl2022therapy}; on the
standardized SEP-28k-E partition per-class F1 varies widely by type (interjection
0.78, prolongation 0.53, block 0.30)~\cite{susac2026stuttmil}. These models encode
linguistic content together with paralinguistic attributes such as speaker
identity and recording condition; how much corpus-specific information a
representation carries is the trade-off our comparison probes.

\subsection{Corpora and Metrics Under Shift}
The three corpora used here (Section~\ref{sec:arch}) differ in speakers,
recording conditions, label definitions, and annotation granularity, so
generalization to unseen conditions is difficult to assess. Language is not among
the differences our evaluation exercises, because the ingested Boli portion is
English-only (Section~\ref{sec:arch}). Bayerl et
al.~\cite{bayerl2023crosscorpus} framed cross-corpus detection as a multi-label
problem, motivated by dysfluent events co-occurring rather than appearing in
isolation, and showed that partitioning changes reported results
substantially~\cite{bayerl2022partitioning}; most of this literature still reports
F1 at a threshold chosen on in-distribution data, the practice whose adequacy we
examine. Our study is methodological: we hold the model set fixed, vary the
evaluation protocol, and report thresholded and threshold-free metrics for the
same predictions.

F1 is the harmonic mean of precision and recall at a chosen operating point, so
its value depends on the threshold as well as on the score ranking. AUROC
summarizes ranking over all thresholds but can look optimistic when positives are
rare, because a large negative class inflates the area; AUPRC weights precision by
the positive class, is the more informative of the two under imbalance, and has
the positive-class prevalence rather than 0.5 as its chance level. Predictive
calibration is a separate property that we do not measure~\cite{guo2017calibration}.
We report thresholded F1, threshold-free AUROC and AUPRC, the AUPRC lift over the
prevalence prior, and the predicted-positive rate (PPR): at a fixed threshold, a
PPR far from the base rate indicates displaced operating-point placement rather
than lost discriminative ability. Because the same thresholds are applied to both
evaluation sets, the difference in PPR between them measures drift directly.

% ── III. CORPORA AND MODELS ──────────────────────────────────────────
\section{Corpora and Models}
\label{sec:arch}

The system is a clip-level classification pipeline that predicts which of the
five dysfluency types are present in a 3-s audio window. All models share one
preprocessing module and are exposed through a common registry API.

\subsection{Source Corpora and Clip Construction}
SEP-28K contributes 32{,}321 clips (28{,}177 podcast clips and a 4{,}144-clip
re-annotation of adult FluencyBank speech); UCLASS contributes 4{,}712
therapy-session clips in SEP-28K format; and Project Boli contributes 54 clips
with word-level annotations and timestamps. Boli is confined to the test split,
so no external speaker is seen in training. The three are merged into 37{,}087
manifest rows, with per-clip interval CSVs recording events as
\texttt{(start\_sec, end\_sec, type)} tuples; clips whose audio is missing or
empty are dropped, leaving 36{,}674 usable clips. The 36{,}620 non-Boli clips are
partitioned 80/10/10 at random under seed 42 into 29{,}296 training, 3{,}662
validation, and 3{,}662 test clips, and all 54 Boli clips are pinned to the test
split, giving a test partition of 3{,}716 rows. Audio is resampled to 16~kHz
mono, DC-removed, peak-normalized, and silence-trimmed.

Evaluation enumerates the audio files in a split directory and keeps those with a
matching label CSV and a non-empty body, so the scored unit is the distinct
recording rather than the manifest row; on the test partition this yields 3{,}715
scored clips, of which 53 are Boli, and the same rule gives the Boli count of 53
used throughout.

\subsubsection{Construction of the Project Boli subset}
\label{sec:boli}
The external subset is derived by a four-step procedure from the released Boli
corpus. First, all \texttt{.txt} files in the \texttt{Transcripts} directory are
enumerated in sorted order; each is a word-level annotation list whose lines carry
a start time, an end time, and a dysfluency code. Second, each transcript is
paired with audio by name (direct filename lookup, then a mapped lookup through a
fixed task table, then a recursive basename search); transcripts with no
resolvable audio are discarded. Third, each transcript is parsed line by line,
lines with fewer than three fields are skipped, and the remaining codes
(\emph{B}, \emph{PR}, \emph{SR}, \emph{WR}, \emph{IN}) are mapped to block,
prolongation, sound repetition, word repetition, and interjection; transcripts
yielding no recognized code are discarded. The 54 survivors each contribute one
manifest row whose five binary labels are set by whether any annotated word
carries the code. Fourth, all 54 rows are pinned to the test split and never seen
in training; scoring enumerates distinct audio filenames, so they resolve to 53
scored recordings, and all external results use those 53 clips. Because the
pairing key is the audio filename, the scored unit is the recording rather than
the annotation row; the reports retain only counts, so we cannot say which filter
removed the single row.

Three properties of the released corpus matter for interpretation. First, each
annotated file is built around a single stuttered word~\cite{boli2025}, so the
external set is enriched for dysfluent content by construction and is not a
sample of naturalistic speech. Second, the released corpus is multilingual (28
participants: 22 Hindi, 2 Telugu, 2 Bengali, 1 Marathi, and 1 Assamese mother
tongue; roughly 2.5~h), but our ingestion path resolves only the English task
codes, so all 54 ingested clips are English. The protocol therefore changes
speakers, recording pipeline, and annotation convention, but not language: it is
a cross-corpus evaluation and not a cross-lingual one, and cross-lingual
robustness remains untested. Third, the retained audio is the first 3~s of each
file (Section~\ref{sec:win}), whereas the Boli label derives from the whole file,
so a recording longer than 3~s may carry a label not audible in the segment the
model receives.

\subsection{Annotation Harmonization}
The three corpora use different conventions, which we map onto five common
binary clip-level labels: prolongation, block, sound repetition, word
repetition, and interjection. SEP-28K labels are adopted as-is, with a clip
positive for a type when at least two annotators marked it
(\texttt{SEP28K\_MIN\_VOTES} $=2$); requiring two votes is a substantive choice,
since accepting any single vote inflates positives several-fold. Because each
UCLASS clip is a 3-s window centered on a stutter event, each present label is
materialized as a weak interval spanning 1.25--1.75~s. Boli tokens
(\emph{B}, \emph{PR}, \emph{SR}, \emph{WR}, \emph{IN}) map through a fixed code
table to the same five classes (Section~\ref{sec:boli}), a clip being positive
when any annotated word carries the class. All models therefore consume identical
five-dimensional binary targets.

\subsection{Preprocessing and the 3-Second Window}
\label{sec:win}
SEP-28K clips are natively 3~s, so a 3-s maximum window matches the predominant
corpus; because the task is clip-level, an event crossing the window boundary
does not alter the target. Longer audio is truncated to its first 3~s and shorter
audio is zero-padded, so every configuration receives an identically shaped input.
Two consequences recur: silence trimming risks suppressing the pause that defines
a block, and for Boli recordings longer than 3~s the label derives from the whole
file while the retained segment is its first 3~s, so a positive label need not be
present in the audio the model receives (Section~\ref{sec:interpret}).

\subsection{Model Configurations}
Seven configurations are evaluated. Trainable parameter counts are 472{,}845{,}450
for arm01 (five independent Wav2Vec2 binary classifiers), 97{,}332{,}362 each for
arm02 and arm03 (Wav2Vec2 multitask, backbone frozen for 3 and 20 epochs), and
341{,}322, 274{,}242, 456{,}906, and 539{,}594 for arm04--arm07 (CNN with pooling,
single-head CNN with pooling, CNN with LSTM, and CNN with transformer). The CNNs
therefore hold 0.3--0.6\% of the parameters of arm01. The evaluation artifacts
report slightly larger CNN counts because they include 708 non-trainable
BatchNorm buffer elements; we quote trainable counts here. Inference latency,
throughput, and peak GPU memory were not measured and are not inferred from these
counts (Section~\ref{sec:compute}).

\subsubsection{Model descriptions.}
arm01 uses five independent binary classifiers on
\texttt{facebook/wav2vec2-base} (94{,}569{,}090 parameters each): raw audio
$\rightarrow$ encoder $\rightarrow$ \texttt{Linear}(768$\rightarrow$256) on the
frame-level hidden states $\rightarrow$ mean pooling over time $\rightarrow$
\texttt{Linear}(256$\rightarrow$2) $\rightarrow$ softmax. The projection precedes
the pooling, the head has no dropout, and no attention mask is passed, so the
temporal reduction is an unmasked mean; fixed 3~s inputs give every clip the same
number of frames. Each classifier is trained independently with focal loss
($\gamma = 2.0$), AdamW (lr $= 3\times10^{-5}$, weight decay $= 0.01$), a 500-step
linear warmup, and mixed precision, with the backbone frozen for 3 epochs and
then trained at one tenth of the head rate ($3\times10^{-6}$). arm02 and arm03
feed five parallel per-class heads from one shared encoder, each an MLP on a
temporal mean pool (\texttt{Linear}(768$\rightarrow$768) $\rightarrow$ Tanh
$\rightarrow$ \texttt{Linear}(768$\rightarrow$2)), minimizing the summed focal
loss: 97.3M parameters in total, only ${\approx}2.8$M more than one independent
classifier. arm03 differs only in freezing the backbone for 20 epochs instead of
3, isolating freezing duration.

arm04--arm07 share a spectrogram backbone and differ in the sequence aggregator.
Input is a log-mel spectrogram with 128 Mel bands, a 2048-point FFT, and hop length
512. The encoder is four convolutional blocks with 32, 64, 128, and 128 channels,
each with a $3\times3$ kernel (padding 1), batch normalization, ReLU, a $(2,1)$
max-pool in all but the last block, and 2-D dropout at rate $0.4$. The three pools
leave a $128\times16\times94$ tensor (channels $\times$ mel bins $\times$ frames).
This is aggregated by global mean pooling plus \texttt{Linear}(128$\rightarrow$128)
(arm04); a single-layer unidirectional LSTM with hidden size 128 over the
frequency-pooled time axis, using the final hidden state (arm06); or a one-layer
transformer (dimension 128, four heads, feed-forward width 512, dropout $0.4$)
followed by temporal mean pooling (arm07). Five per-class two-layer MLP heads
(\texttt{Linear}(128$\rightarrow$128) $\rightarrow$ Tanh $\rightarrow$
\texttt{Linear}(128$\rightarrow$2)) produce the logits. The CNNs are trained from
scratch, with no backbone freezing or pre-training, using binary cross-entropy on
the positive logit with a per-class positive weight equal to the
negative-to-positive ratio of the training pool, summed over classes. The families
thus differ in training objective as well as in representation, which we note as a
confound. arm05 is the single-head variant of the pool model: the head block is
built once per class rather than five times, and one model is trained per
dysfluency type (five in total), mirroring the five-binary Wav2Vec2 design at a
much smaller scale.

% ── IV. EXPERIMENTAL SETUP ───────────────────────────────────────────
\section{Experimental Setup}
\label{sec:setup}

\subsection{Training Protocol}
All models are trained on a single NVIDIA GPU with seed $= 42$, 20 epochs and
early stopping (patience $= 5$), mixed precision, batch size 8 (binary
classifiers) or 16 (multitask and CNN), and 3~s audio at 16~kHz. Training reads
the training directory only: it derives its own held-out portion, a 20\%
stratified subsample (5{,}857 clips, seed 42) used for early stopping and
threshold tuning, so parameters are fitted on 23{,}439 clips and the prepared
validation split is never opened. No gradient clipping or input augmentation is
applied; an earlier gradient-clipping configuration was abandoned because, on a
94.6M-parameter model with a randomly initialized head, the clip dominated the
head's early gradient signal. Each configuration is trained once, so every
reported value is a single-run point estimate.

\subsection{Threshold Selection}
Thresholds are tuned on the same 20\% subsample by sweeping 17 fixed values from
0.1 to 0.9 in steps of 0.05 and, per class, selecting the value maximizing F1
rather than Youden's J. Because that subsample is also the early-stopping set,
the tuned thresholds are mildly optimistic; we treat this as a limitation rather
than as test-set leakage, since the test split is never used for any selection.
arm02 uses thresholds of 0.55, 0.35, 0.40, 0.35, and 0.45 for
prolongation, block, sound repetition, word repetition, and interjection. No
test-set or cross-corpus threshold fitting is performed: the same thresholds are
applied unchanged to the external set, the transfer that
Sections~\ref{sec:results} and~\ref{sec:discussion} examine.

\subsection{Reported Metrics}
\label{sec:metrics}
For each configuration, protocol, and class we report precision, recall, F1,
AUROC, and AUPRC, plus macro averages (the unweighted mean over the five classes),
which keep classes of differing prevalence from being dominated by the most
frequent. AUPRC is reported as a lift over its chance level, the positive-class
prevalence. Precision and recall are omitted from the main tables, having fixed
the threshold in advance, but are recoverable from F1 and the PPR with the
confusion counts.

\subsection{Evaluation Protocols}
\subsubsection{In-distribution, clip-level (Test).}
The test partition holds 3{,}716 manifest rows, of which 3{,}662 come from
SEP-28K and UCLASS and 54 from Boli. The partition operates on clips, not
speakers, so speakers may appear in both training and test for the SEP-28K and
UCLASS portions; these numbers are an upper-bound, not a speaker-independent,
benchmark. Results are reported over the 3{,}715 scored clips, of which the 53
scored Boli recordings (${\approx}1.4\%$) are out-of-domain. The macro base rate
is 0.143.

\subsubsection{Speaker-independent, cross-corpus (Boli).}
The 53 scored Boli recordings are unseen during training, with different
speakers, recording conditions, and label-construction procedure
(Section~\ref{sec:arch}). Per-class support is small: PR 16, BL 35, SR 34, WR 14,
IN 9, giving a macro base rate of 0.408. We treat all external conclusions as
exploratory.

Two trivial baselines accompany the models: the \emph{all-positive} classifier,
which labels every clip positive for every class, and the \emph{random-at-rate}
classifier, which labels clips positive with probability equal to the class base
rate. All-positive is the more informative under heavy imbalance: it attains
recall 1 by construction and its F1 is determined entirely by the base rate, so a
model that does not exceed it has not demonstrated usable discrimination; the
expected macro F1 of random-at-rate equals the macro base rate.

\subsection{Statistical Uncertainty}
\label{sec:uncert}
Uncertainty is quantified by nonparametric bootstrap resampling of clips: for each
reported quantity we resample the evaluated clips with replacement 10{,}000 times,
recompute the metric, and take the 2.5th and 97.5th percentiles. Resampling is
unstratified, so it propagates variability in class prevalence as well as in the
decisions; on the 53-clip external set the former dominates for the trivial
baselines. F1 at a fixed threshold and the predicted-positive rate are functions
of the confusion counts alone, so their bootstrap operates exactly on the stored
confusion matrices. These intervals capture clip-level sampling variability only:
each configuration was trained once, so they exclude training seed, and
overlapping intervals should not be read as a real difference in quality. AUROC
and AUPRC are point estimates without intervals, because both depend on the full
per-clip score vector, which was not retained; we therefore treat the
threshold-free ordering as directional evidence.

% ── V. RESULTS ───────────────────────────────────────────────────────
\section{Results}
\label{sec:results}

\subsection{In-Distribution Results (RQ1)}
\label{sec:res-indist}
Table~\ref{tab:macro} reports macro F1 for all seven configurations under fixed
(0.5) and tuned thresholds on both evaluation sets, with the two trivial
baselines. Parenthesized intervals are 95\% bootstrap confidence intervals
(Section~\ref{sec:uncert}).
\begin{table}[!htbp]
\caption{Macro F1 (unweighted mean over five classes) with 95\% bootstrap
intervals on tuned columns. Test: clip-level in-distribution split (3{,}715
clips, speakers shared with training); Boli: speaker-independent external split
(53 recordings). Tuned thresholds are unchanged on Boli; baselines are untuned.}
\label{tab:macro}
\centering
\footnotesize
\setlength{\tabcolsep}{2pt}
\renewcommand{\arraystretch}{0.9}
\begin{tabular}{lrrrr}
\hline
Arm & Test F1 (0.5) & Test F1 (tuned) & Boli F1 (0.5) & Boli F1 (tuned)\\
\hline
arm01\_5x\_w2v2 & 0.573 & 0.570 \ci{0.554--0.587} & 0.254 & 0.292 \ci{0.169--0.391}\\
arm02\_mt\_frz3 & 0.542 & 0.556 \ci{0.539--0.572} & 0.162 & 0.238 \ci{0.127--0.332}\\
arm03\_mt\_frz20 & 0.141 & 0.378 \ci{0.362--0.394} & 0.046 & 0.214 \ci{0.119--0.298}\\
arm04\_cnn\_pool & 0.149 & 0.256 \ci{0.245--0.266} & 0.351 & 0.332 \ci{0.254--0.398}\\
arm05\_cnn\_single & 0.212 & 0.262 \ci{0.251--0.272} & 0.450 & 0.476 \ci{0.389--0.547}\\
arm06\_cnn\_lstm & 0.242 & 0.249 \ci{0.240--0.257} & 0.451 & 0.508 \ci{0.440--0.566}\\
arm07\_cnn\_tf & 0.269 & 0.271 \ci{0.259--0.283} & 0.421 & 0.345 \ci{0.260--0.419}\\
\hline
All-positive & 0.247 & 0.247 \ci{0.239--0.254} & 0.550 & 0.550 \ci{0.490--0.604}\\
Random@rate & 0.143 & 0.143 & 0.408 & 0.408\\
\hline
\end{tabular}
\end{table}

Under the clip-level protocol, arm01 attains the highest measured test macro F1
at tuned thresholds and arm02 follows closely: 0.570 (0.554--0.587) against 0.556
(0.539--0.572), with substantially overlapping intervals on 3{,}715 clips. Since
the intervals describe clip sampling alone and each configuration was trained
once, we report the pair as indistinguishable and do not treat arm01 as the
better model. Both exceed every CNN by a wide margin, their intervals lying far
above the CNN range. The CNNs reach macro F1 of 0.25--0.27 at tuned thresholds,
above the random-at-rate baseline (0.143) but only marginally above the
all-positive baseline (0.247); for arm06 the two intervals overlap. Freezing
duration is the one architectural factor with a large and unambiguous effect:
arm03, which freezes the backbone for 20 epochs rather than 3, falls to 0.378
(0.362--0.394), a relative reduction of ${\approx}32\%$ from arm02 with
non-overlapping intervals, evidence that task adaptation of the encoder is
necessary here. Its fixed-0.5 score of 0.141 also falls below random-at-rate,
suggesting poorly distributed rather than merely poorly ranked scores.

\subsection{Cross-Corpus Results at Tuned Thresholds (RQ2)}
\label{sec:res-boli-tuned}
On the 53-clip external set, the all-positive baseline attains macro F1 0.550
(0.490--0.604), which no configuration exceeds; the best model, arm06, reaches
0.508 (0.440--0.566). The intervals overlap, so this compares point estimates
rather than a separated difference, and the CNN ordering is likewise unresolved:
the four CNN intervals span 0.254 to 0.566 and overlap extensively. The one clear
separation is between the best CNN and the two task-adapted Wav2Vec2
configurations, whose intervals (0.169--0.391 and 0.127--0.332) lie below that of
arm06. Table~\ref{tab:boli_perclass} gives the per-class breakdown.

\begin{table}[!htbp]
\caption{Per-class F1 on the external set (53 recordings) at the
validation-tuned thresholds, transferred unchanged. All-positive labels every
clip positive, giving per-class F1 $2s/(n+s)$ for support $s$ and $n=53$; no
configuration exceeds its macro F1. Macro intervals are 95\% bootstrap intervals;
per-class intervals are omitted because with support 9--35 they are too wide to
inform the comparison.}
\label{tab:boli_perclass}
\centering
\small
\renewcommand{\arraystretch}{0.9}
\begin{tabular}{lrrrrrr}
\hline
Arm & PR & BL & SR & WR & IN & Macro\\
\hline
arm01\_5x\_w2v2 & 0.400 & 0.250 & 0.256 & 0.125 & 0.429 & 0.292 \ci{0.169--0.391}\\
arm02\_mt\_frz3 & 0.200 & 0.154 & 0.256 & 0.182 & 0.400 & 0.238 \ci{0.127--0.332}\\
arm03\_mt\_frz20 & 0.200 & 0.233 & 0.182 & 0.182 & 0.273 & 0.214 \ci{0.119--0.298}\\
arm04\_cnn\_pool & 0.308 & 0.491 & 0.182 & 0.387 & 0.290 & 0.332 \ci{0.254--0.398}\\
arm05\_cnn\_single & 0.326 & 0.765 & 0.508 & 0.431 & 0.348 & 0.476 \ci{0.389--0.547}\\
arm06\_cnn\_lstm & 0.492 & 0.795 & 0.542 & 0.418 & 0.290 & 0.508 \ci{0.440--0.566}\\
arm07\_cnn\_tf & 0.316 & 0.320 & 0.490 & 0.267 & 0.333 & 0.345 \ci{0.260--0.419}\\
\hline
All-positive & 0.464 & 0.795 & 0.782 & 0.418 & 0.290 & 0.550 \ci{0.490--0.604}\\
\hline
\end{tabular}
\end{table}

The CNN rows appear to show that compact spectrogram models ``generalize best'';
that reading is incorrect, and Table~\ref{tab:boli_drift} shows why. The CNNs
reach their higher F1 on block, word repetition, and interjection by labelling
nearly every external clip positive: arm06 applies the positive label to all 53
clips for each of those three classes, so its per-class F1 approaches, and for
interjection exactly equals, the value a predictor with no learned discrimination
attains at the same support. Its macro PPR is 0.879, more than twice the 0.408
base rate, while the Wav2Vec2 configurations move the opposite way at PPR
0.079--0.102, leaving their F1 limited by near-zero recall. Both behaviours follow
from where a fixed threshold lands after the score distribution shifts; neither
is evidence about ranking quality. The finding is negative: at thresholds never
fitted to the external set, no configuration demonstrates usable dysfluency
discrimination, and the CNN advantage under thresholded F1 is consistent with
operating-point drift rather than superior ranking.

\subsection{Threshold-Free Cross-Corpus Results and Operating-Point Drift (RQ3, RQ4)}
\label{sec:res-thresholdfree}
Because F1 conflates ranking with threshold placement, we also report
threshold-free quantities for the same predictions
(Table~\ref{tab:boli_drift}), together with the predicted-positive rate measured
on the in-distribution split under the same thresholds. AUPRC lift is AUPRC
divided by the macro base rate of 0.408, the AUPRC of a random predictor.

\begin{table}[!htbp]
\caption{Threshold-free macro metrics and predicted-positive rate (PPR) from the
same predictions as Table~\ref{tab:boli_perclass}. AUROC and AUPRC are point
estimates; AUPRC lift is relative to the 0.408 base rate. PPR is the mean
fraction labelled positive at the tuned thresholds and carries a 95\% bootstrap
interval. The PPR columns use identical thresholds on both sets, so their
difference measures drift directly. Base-rate rows give the in-distribution
(0.143) and external (0.408) priors.}
\label{tab:boli_drift}
\centering
\small
\renewcommand{\arraystretch}{0.9}
\setlength{\tabcolsep}{2pt}
\begin{tabular}{lrrrrr}
\hline
Arm & AUROC & AUPRC & AUPRC lift & PPR (Test) & PPR (Boli)\\
\hline
arm01\_5x\_w2v2 & 0.620 & 0.546 & 1.34 & 0.134 & 0.079 \ci{0.049--0.113}\\
arm02\_mt\_frz3 & 0.621 & 0.557 & 1.37 & 0.159 & 0.102 \ci{0.068--0.140}\\
arm03\_mt\_frz20 & 0.538 & 0.443 & 1.09 & 0.195 & 0.204 \ci{0.158--0.253}\\
arm04\_cnn\_pool & 0.473 & 0.412 & 1.01 & 0.572 & 0.574 \ci{0.528--0.619}\\
arm05\_cnn\_single & 0.497 & 0.445 & 1.09 & 0.583 & 0.615 \ci{0.566--0.664}\\
arm06\_cnn\_lstm & 0.522 & 0.419 & 1.03 & 0.845 & 0.879 \ci{0.849--0.909}\\
arm07\_cnn\_tf & 0.500 & 0.412 & 1.01 & 0.365 & 0.438 \ci{0.381--0.494}\\
\hline
Base rate (random) & 0.500 & 0.408 & 1.00 & 0.143 & 0.408\\
\hline
\end{tabular}
\end{table}

The threshold-free view inverts the thresholded ordering once more. arm01 and
arm02 rank external clips best, with macro AUROC 0.620--0.621 and macro AUPRC
0.546--0.557, a lift of ${\approx}1.3$--$1.4\times$ over the prevalence prior;
arm03 sits near chance; and the CNNs are at or barely above chance (AUROC
0.473--0.522, AUPRC 0.412--0.445). The advantage is modest: these are point
estimates, and on 53 clips with at most 35 positives per class and 9 in the worst
case, a macro AUROC of 0.62 against 0.52 is well within the range that resampling
could reduce to nothing. We therefore state the result directionally: the
threshold-free metrics place the task-adapted Wav2Vec2 configurations above the
CNNs, and no configuration comes near reliable cross-corpus detection.

The same table answers RQ3: both PPR columns use identical thresholds, so the
comparison is a measurement rather than an inference. For six of seven
configurations the 95\% interval for external macro PPR excludes 0.408, the
exception being arm07 (0.381--0.494); the Wav2Vec2 configurations sit below
(0.079--0.204) and the CNNs above (0.438--0.879). The CNN operating points were
already displaced in-distribution: at base rate 0.143 they label 0.365--0.845 of
clips positive (2.6--5.9 times prevalence) while reaching only 0.249--0.271 macro
F1. Their apparent external F1 gain is thus not produced by the shift: the
over-positive operating point that cripples in-distribution F1 becomes
advantageous externally only because the Boli base rate is nearly three times
higher. In relative terms their operating point is \emph{less} displaced
externally (arm06: 5.9 times the base rate on test versus 2.2 times on Boli), the
opposite of the drift a generalization failure would produce. The task-adapted
Wav2Vec2 configurations move the other way under the same thresholds, from 0.134
and 0.159 in-distribution to 0.079 and 0.102 externally, toward all-negative.
Neither behaviour is about ranking quality, because AUROC and AUPRC do not move
with the threshold and order the two families the other way round.


\subsection{Per-Class Analysis}
Table~\ref{tab:perclass} gives per-class metrics for arm02 on the
clip-level test split at its tuned thresholds.

\begin{table}[!htbp]
\caption{Per-class metrics for the multitask Wav2Vec2 model (arm02) on the
in-distribution test split (3{,}715 clips) at validation-tuned thresholds 0.55,
0.35, 0.40, 0.35, and 0.45 for prolongation, block, sound repetition, word
repetition, and interjection. Base rate is support divided by 3{,}715, the chance
level of AUPRC. F1 carries a 95\% bootstrap interval.}
\label{tab:perclass}
\centering
\footnotesize
\setlength{\tabcolsep}{3pt}
\renewcommand{\arraystretch}{0.9}
\begin{tabular}{lrrrrr}
\hline
Class & F1 (tuned) & AUROC & AUPRC & Support & Base rate\\
\hline
Prolongation & 0.457 \ci{0.407--0.503} & 0.845 & 0.491 & 400 & 0.108\\
Block & 0.429 \ci{0.397--0.460} & 0.768 & 0.444 & 553 & 0.149\\
Sound Rep. & 0.621 \ci{0.584--0.656} & 0.878 & 0.624 & 494 & 0.133\\
Word Rep. & 0.499 \ci{0.456--0.541} & 0.855 & 0.550 & 366 & 0.099\\
Interjection & 0.773 \ci{0.750--0.795} & 0.935 & 0.848 & 834 & 0.225\\
\hline
\textbf{Macro} & \textbf{0.556 \ci{0.540--0.572}} & \textbf{0.856} & \textbf{0.592} & --- & \textbf{0.143}\\
\hline
\end{tabular}
\end{table}

Interjection is the easiest class, at F1 0.773 (0.750--0.795) with AUROC 0.935,
and block the hardest, at F1 0.429 (0.397--0.460) with AUROC 0.768; the intervals
are well separated. Read against their chance levels, the AUPRC values in
Table~\ref{tab:perclass} are 2.99--5.59 times the base rate, so the ranking signal
is real in-distribution even where thresholded F1 is low. One reading is that
interjections are stereotyped fillers whose acoustic form differs from fluent
speech, whereas blocks are pauses close to ordinary hesitation that depend on
prosodic and contextual information a 3-s window and a mean-pooled representation
may not preserve; the mechanisms are not separable here, and the block result
fits the window and trimming artifacts (Section~\ref{sec:arch}) equally well.

For every class AUROC exceeds tuned F1, most markedly for prolongation (0.845
against 0.457): the models order dysfluent clips better than the tuned thresholds
convert that ordering into decisions. Whether this reflects miscalibration or a
mismatch between F1-optimal and accuracy-optimal operating points cannot be
settled here, because we do not measure calibration. The arm01 block classifier
produces the most false negatives (343) and the most false positives relative to
its support (284).

\subsection{Computational Cost and Relation to Published Results}
\label{sec:compute}
Parameter count, storage size, computational cost, and runtime cost are distinct:
the first two follow from the weights and a stated precision, the last must be
measured. We therefore report parameter counts and their fp32 storage in
Section~\ref{sec:arch} but not latency, throughput, or peak GPU memory; parameter
counts give fp32 storage of 1.89\,GB for arm01 against 1.1--2.2\,MB for the CNNs.
Within the Wav2Vec2 family we would select the multitask configuration
operationally: 97.3M parameters against 472.8M because the encoder is shared, and
one forward pass per clip instead of five. It does not measurably beat arm01. Our
best in-distribution macro F1, 0.570 (five-binary) and 0.556 (multitask), falls
within the 0.56--0.63 reported for Wav2Vec2 on FluencyBank~\cite{bayerl2022multi},
and our per-class ordering matches the 0.30 (block) to 0.78 (interjection) span on
SEP-28k-E~\cite{susac2026stuttmil}; this is an ordering comparison, not a numerical
one, since our split shares speakers and the corpora, label definitions, and
threshold strategies differ~\cite{bayerl2022partitioning}.

% ── VI. DISCUSSION ───────────────────────────────────────────────────
\section{Discussion}
\label{sec:discussion}

\subsection{What the In-Distribution Results Show}
The Wav2Vec2 configurations capture in-distribution signal that the compact CNNs,
at roughly 0.3--0.6\% of the parameter count, do not. The most likely explanation
is representational capacity, but capacity, pretraining, and objective are not
separable in our design, since the families also differ in loss and we vary the
architecture as a whole; we therefore do not attribute the gap to a single
factor. The one isolable factor, freezing duration, acts in the opposite
direction and is unambiguous: withholding task adaptation from the same encoder
costs ${\approx}32\%$ of tuned macro F1, an effect larger than any difference
among the CNNs.

\subsection{Why the External Results Are Hard to Interpret}
\label{sec:interpret}
Three features complicate readings of Tables~\ref{tab:boli_perclass}
and~\ref{tab:boli_drift}: 53 recordings, per-class support as low as 9, and labels
constructed unlike the training corpora. The intervals quantify the first two;
the third is most consequential. The strongest model's external macro F1, 0.508
(0.440--0.566), overlaps the all-positive baseline's 0.550 (0.490--0.604), so the
negative result concerns point estimates on a set too small to separate them. The
comparison the data support cleanly is the operating-point one: PPR intervals
exclude the prevailing base rate for six of seven configurations, and the same
thresholds in-distribution show the CNN operating points were already displaced.
That measurement, not the F1 ranking, is the substantive content.

Because Boli files are built around an annotated stutter event~\cite{boli2025}, a
positive label reflects a curated event rather than sampling continuous speech and
applying a majority vote, as with SEP-28K and UCLASS. The all-positive baseline
consequently scores 0.550 here against 0.247 on test, a gap at least partly due to
clip selection rather than any model. A second complication is the label--window
mismatch, which depresses recall unrelated to representation family; we cannot
determine how often it occurs because durations were not retained.

\subsection{Proposed Explanations}
\label{sec:hypotheses}
Two explanations the present experiments do not distinguish are stated as
proposals rather than findings.

\emph{Operating-point drift.} The score distributions of all seven configurations
are displaced, so thresholds fixed in-distribution land in different ROC regions
and F1 tracks the predicted-positive rate and base rate rather than
discrimination. Direct evidence is the paired PPR columns of
Table~\ref{tab:boli_drift}, where the families sit on opposite sides of the base
rate and the ordering flips when the threshold is removed; because both use the
same thresholds, displacement is measured rather than assumed. Why it occurs
would require the score distributions themselves; the F1 versus AUROC/AUPRC
disagreement is its expected signature.

\emph{Corpus-specificity.} Higher-capacity Wav2Vec2 representations may encode
corpus-specific characteristics (speakers, accents, conditions, annotation
conventions) that transfer weakly, while coarse spectral patterns transfer more
evenly. The threshold-free ordering is consistent, but so are the label--window
mismatch and label provenance; no measurement separates them. A probe classifier
on frozen embeddings would be the direct test; we have not run one.

\subsection{Threats to Validity}
\label{sec:limitations}
\paragraph{Statistical power and single-run training.} This is the most
consequential threat, with two sources. Each configuration was trained once from
seed 42, so the arm01/arm02 difference of under 0.02 macro F1 is not interpretable
even with overlapping intervals, and the arm03 degradation carries no training
uncertainty: the intervals in Tables~\ref{tab:macro}--\ref{tab:perclass} describe
clip sampling only. The external set's 53 recordings with support 9--35 also
explain why external F1 comparisons are point estimates and threshold-free
metrics have no intervals. We use ``better'' and ``outperforms'' only where a
difference exceeds plausible seed and sampling variation.

\paragraph{Uncontrolled properties of the external set, and untested
explanations.} The external labels are built differently from training labels,
the clips are English-only and do not test language shift, and truncation to the
first 3~s may decouple labels from retained audio. The window and trimming policy
is fixed across configurations and protocols, so it does not confound comparisons,
but it is an uncontrolled absolute factor and a plausible contributor to block
difficulty. Tuned F1 depends on validation-optimal thresholds, and without
calibration we cannot distinguish a miscalibrated head from a mismatched operating
point~\cite{guo2017calibration}; the two accounts remain proposals. Capacity,
pretraining, and objective are not separated, latency was not measured, and one
external corpus cannot distinguish a general transfer failure from a corpus
property.

\section{Practical Implications and Conclusion}
\label{sec:conclusion}
Model selection for this task cannot rest on in-distribution F1, because that
split shares speakers with training and the cross-corpus ordering is dominated by
where a fixed threshold lands. Reported results should therefore include a
speaker-disjoint or cross-corpus split, threshold-free metrics alongside F1, an
all-positive baseline, and the predicted-positive rate at the reported threshold.

On RQ1, Wav2Vec2 dominates in-distribution at macro F1 0.55--0.57 against
0.25--0.27 for the CNNs and 0.247 for an all-positive baseline, with the
five-binary and multitask variants indistinguishable and excessive freezing
costing ${\approx}32\%$ of tuned macro F1. On RQ2 the cross-corpus F1 ordering
inverts, with the CNNs appearing to lead, and on RQ3 that inversion is consistent
with operating-point drift rather than generalization: at identical thresholds the
CNNs label 37--85\% and 44--88\% of in-distribution and external clips positive
against base rates of 0.143 and 0.408, while the Wav2Vec2 configurations fall to
8--10\% externally, and no configuration exceeds the all-positive baseline of
0.550. On RQ4 threshold-free metrics reverse the ordering again, placing the
task-adapted Wav2Vec2 configurations above the CNNs at macro AUROC 0.620--0.621
against 0.473--0.522. The conclusions we can draw are narrow: F1 alone cannot
support cross-corpus claims for this task, and interjection is easier to detect
than block under both protocols. We cannot claim cross-corpus deployment, because
no configuration exceeded a trivial baseline on the external set and at least
three confounding properties remain (Section~\ref{sec:arch}). The cause of the
score drift and the modest threshold-free advantage of Wav2Vec2 remain unresolved
(Section~\ref{sec:limitations}).

% ---- Bibliography ----
%
% BibTeX users should specify bibliography style 'splncs04'.
% References will then be sorted and formatted in the correct style.
%
\bibliographystyle{splncs04}
\bibliography{references}

% ── CONFLICT OF INTEREST ─────────────────────────────────────────────
\begin{credits}
\subsubsection{\discintname}
The authors declare no competing interests relevant to this article.
\end{credits}
%
\end{document}
